\documentclass[11pt]{article}
\usepackage[a4paper,margin=1in]{geometry}
\usepackage{fontspec}
\usepackage[boldfont=false]{xeCJK}
\setCJKmainfont{FandolSong-Regular.otf}
\usepackage{amsmath}
\usepackage{amssymb}
\usepackage{array}
\usepackage{booktabs}
\usepackage{graphicx}
\usepackage{adjustbox}
\usepackage{enumitem}
\setlist[itemize]{topsep=2pt,partopsep=0pt,parsep=0pt,itemsep=2pt,leftmargin=1.4em}
\setlist[enumerate]{topsep=2pt,partopsep=0pt,parsep=0pt,itemsep=2pt,leftmargin=1.6em}
\usepackage{parskip}
\usepackage{fvextra}
\fvset{breaklines=true,breakanywhere=true,fontsize=\small}
\setlength{\parindent}{0pt}

\begin{document}

\begin{center}
  {\LARGE\bfseries The mitotic cell cycle}\\[0.35em]
  {\large A Level Biology}
\end{center}
\vspace{0.6em}

\section*{The structure of a chromosome}

A \textbf{chromosome} 染色体 is one very long molecule of DNA wound tightly around special \textbf{proteins} 蛋白质 called \textbf{histones} 组蛋白. Winding the DNA like this lets a huge length fit inside the nucleus and keeps it tidy.

Before a cell divides, its DNA is copied (this copying is called \textbf{replication} 复制). After copying, each chromosome is made of two identical copies joined together. These two copies are the \textbf{sister chromatids} 姐妹染色单体, and they are held together at a point called the \textbf{centromere} 着丝粒. The tips of each chromosome are capped by \textbf{telomeres} 端粒, which protect the ends.

\section*{Why mitosis matters}

\textbf{Mitosis} 有丝分裂 is a type of nuclear division that makes two \textbf{daughter cells} 子细胞 that are \textbf{genetically identical} — they carry exactly the same \textbf{genes} 基因 as the parent cell and as each other.

This matters for:

\begin{itemize}
\item \textbf{growth} of \textbf{multicellular} 多细胞 \textbf{organisms} 生物体 (making more cells).

\item \textbf{replacement} of damaged or dead cells.

\item \textbf{repair} of \textbf{tissues} 组织 by making new cells.

\item \textbf{asexual reproduction} 无性生殖 (one parent makes identical offspring).

\end{itemize}

\section*{The mitotic cell cycle}

The \textbf{cell cycle} 细胞周期 is the full life of a cell from one division to the next. It has three parts:

\begin{enumerate}
\item \textbf{interphase} 间期 — the longest part. The cell grows in the \textbf{G₁} phase, copies its DNA in the \textbf{S} phase (replication), and grows again and prepares to divide in the \textbf{G₂} phase.

\item \textbf{mitosis} — the nucleus divides into two identical nuclei.

\item \textbf{cytokinesis} 胞质分裂 — the rest of the cell splits, giving two separate daughter cells.

\end{enumerate}

\section*{Telomeres and the ends of chromosomes}

When DNA is replicated, the copying cannot reach the very end of the molecule, so a little is lost each time. Telomeres are short, repeated lengths of DNA at the ends that carry \textbf{no genes}. Because the telomeres are shortened instead, no important genes are lost during replication.

\section*{Stem cells}

A \textbf{stem cell} 干细胞 is an unspecialised cell that can keep dividing by mitosis and can \textbf{differentiate} 分化 (change) into different specialised cell types. Stem cells are the source of new cells for replacing lost cells and repairing tissues.

\section*{Uncontrolled division and tumours}

The cell cycle is normally tightly controlled, so cells divide only when needed. If this control is lost, a cell may divide again and again without stopping. This uncontrolled division produces a lump of cells called a \textbf{tumour} 肿瘤.

\section*{The stages of mitosis}

Mitosis runs through four stages. You should be able to recognise them in photomicrographs and slides.

\begin{center}
\adjustbox{max width=\linewidth}{%
\begin{tabular}{ll}
\toprule
\textbf{Stage} & \textbf{What happens} \\
\midrule
\textbf{prophase} 前期 & chromosomes coil up and become visible as two sister chromatids; the \textbf{nuclear envelope} 核膜 breaks down; a \textbf{spindle} 纺锤体 of fibres forms across the cell \\
\textbf{metaphase} 中期 & chromosomes line up along the middle (the \textbf{equator} 赤道); spindle fibres attach to each centromere \\
\textbf{anaphase} 后期 & the centromeres split; the sister chromatids are pulled to opposite ends (poles) of the cell \\
\textbf{telophase} 末期 & a set of chromosomes reaches each pole; a new nuclear envelope forms around each set, making two nuclei \\
\bottomrule
\end{tabular}}
\end{center}

Cytokinesis then follows. In an animal cell the cell surface membrane pinches inwards to split the cell; in a plant cell a new wall forms across the middle. The result is two genetically identical daughter cells.


\end{document}
